% Ad Familiares 13.10 (ad-familiares-13-10) --- English translation
% Translated by Claude (Anthropic) from the Latin (Perseus).
% Style guide: STYLE.md.

\ciceroLetterOpener{CICERO BRVTO S.}

\ciceroSection{1} When your quaestor, M. Terentius, was setting out to join you, I did not think he needed any recommendation from me; I judged he was sufficiently recommended to you already by the ancestral custom itself, which, as you well know, has held this bond of the quaestorship to be next in closeness to the tie of one's own children. But since he had got it into his head that a carefully written letter from me about him would carry the greatest weight with you, and since he kept pressing me to write with the utmost care, I preferred to do what a friend of mine considered so much to his own interest.

\ciceroSection{2} So that you may understand why I owe him this: as soon as M. Terentius came of age and took his place in the Forum, he attached himself to my friendship; then, once he had matured, two further reasons came in to increase my goodwill toward him --- first that he was active in this study of ours, in which I still take the greatest delight, and with intellect, as you know, and not without industry; second that he early threw himself into the partnerships of the public revenues. That last I could have wished he had not done; for he sustained very heavy losses --- but the common cause of an order I hold dear in the highest degree made our friendship stronger. Then, after appearing on both sides of the bench with the best repute alike for honesty and for credit, even before this revolution in public affairs, he stood for office and accounted the most honourable office the proper reward of his labour.

\ciceroSection{3} In these recent times, what is more, he set out for Caesar from Brundisium at my instance, carrying letters and instructions from me; and in that errand I had ample proof both of his affection in undertaking the business and of his faithfulness in carrying it through and reporting back. I feel that, since I had meant to speak separately about his integrity and his character, by first setting out for you the reason why I love him so warmly I have, in the very setting-out of the cause, said enough about his integrity too. But still, separately I promise --- and take it on myself --- that he will be both a pleasure and a use to you: you will find him a man of moderation and of sense, perfectly free from any cupidity, and besides this a man of great capacity for work and the highest industry.

\ciceroSection{4} Nor ought I to be making promises in matters which, once you have come to know him well, are yours to judge; but still, in all newly formed connections it matters what the first approach is, and by what recommendation the door, so to speak, of friendship is opened. That is what I have wanted to do by this letter; though the quaestorship bond itself ought to have accomplished it on its own, the thing is still no weaker for having this added. See to it then, if you value me as much as Varro thinks you do and as I myself feel you do, that I learn as soon as possible that this recommendation of mine has brought him as much advantage as he himself hoped, and as I did not doubt it would.
